\section{Bài toán}

\begin{frame}{Bài toán}{Hỏi--đáp có trích dẫn trên \NumCorpusArticles{} bài báo CNN}
  \vspace{0.5em}
  \begin{center}
    \small
    \begin{tabular}{c c c c c}
      \textbf{Câu hỏi} & $\rightarrow$ &
      \fbox{\parbox{3.4cm}{\centering tìm đoạn văn\\[-0.2em]{\scriptsize\color{brandaccent}\textbf{PHASE 1}}}} &
      $\rightarrow$ &
      \fbox{\parbox{3.4cm}{\centering viết đáp án\\[-0.2em]{\scriptsize\color{brandaccent}\textbf{PHASE 2}}}} \\
    \end{tabular}
    \\[0.8em]
    $\rightarrow$\quad đáp án ngắn \textbf{kèm số trích dẫn} \texttt{[n]}
  \end{center}

  \vspace{0.8em}
  Người dùng hỏi một câu. Hệ thống phải tìm đúng đoạn văn chứa câu trả lời
  trong \NumCorpusChunks{} đoạn, rồi để LLM viết đáp án \emph{chỉ dựa trên} các
  đoạn đó và chỉ ra nó lấy từ đoạn nào.

  \takeaway{Hai việc, hai kiểu hỏng, hai loại thước đo --- nên tách làm hai phase.}
\end{frame}

\begin{frame}{Vì sao phải tách hai phase}
  \vspace{0.3em}
  \small
  \begin{tabularx}{\textwidth}{@{}p{2.6cm} L L@{}}
    \toprule
     & \textbf{Phase 1 --- lấy đoạn văn} & \textbf{Phase 2 --- hỏi LLM} \\
    \midrule
    Hỏng kiểu gì & Không tìm ra đoạn chứa đáp án & Tìm ra rồi nhưng trả lời sai, hoặc bịa \\
    Đo bằng gì & Hit@k, nDCG@5, MRR@5 & Answer Correctness, Faithfulness, Citation \\
    Chi phí & GPU, \textbf{không gọi API} & \textbf{Tốn tiền API} mỗi lần chạy \\
    Tái lập & Có (chỉ mục cố định) & Chỉ trong giới hạn nhiễu của LLM \\
    \bottomrule
  \end{tabularx}

  \vspace{0.8em}
  \textbf{Thứ tự bắt buộc:} khóa Phase 1 trước. Nếu vừa đổi cách lấy đoạn vừa
  đổi cách hỏi, không cách nào biết được cải thiện đến từ đâu.

  \takeaway{Phase 2 chạy trên một vết truy xuất \emph{đóng băng}: mọi khác biệt quy hết về tầng sinh.}
\end{frame}
